\documentclass[11pt]{article}
\usepackage[margin=1in]{geometry}
\usepackage{amsmath,amssymb,amsthm,graphicx,microtype,enumitem}
\usepackage[hidelinks]{hyperref}

\newtheorem{theorem}{Theorem}
\newtheorem{lemma}[theorem]{Lemma}
\newtheorem{proposition}[theorem]{Proposition}
\theoremstyle{remark}
\newtheorem{remark}[theorem]{Remark}
\newcommand{\Md}{M_d(\mathbb C)}
\newcommand{\Tr}{\operatorname{Tr}}
\newcommand{\norm}[1]{\left\lVert#1\right\rVert}

\title{Unique extension without a characteristic sequence:\\
a sharp full-matrix dimension threshold}
\author{Candidate full counterexample to Question 1 of arXiv:1709.01649}
\date{13 August 2026}

\begin{document}
\maketitle

\begin{center}
\fbox{\parbox{0.91\textwidth}{\textbf{Status.}
Candidate full negative answer, likely valid.  The question fails for an
explicit pure state and a generating codimension-one operator system in
$M_3(\mathbb C)$.  More generally, it fails in every $M_d$ with $d\geq3$,
whereas it holds for every operator system in $M_1$ and $M_2$.  Thus dimension
three is the sharp threshold among full matrix algebras.}}
\end{center}

\section{The source question}

Let $A$ be a unital $C^*$-algebra and $S\subset A$ an operator system.  A
sequence $(\Delta_n)$ in $A$ is \emph{characteristic} for a state $\psi$ if
\begin{enumerate}[label=(\alph*)]
\item $\norm{\Delta_n}=1$ for every $n$;
\item $\psi(\Delta_n)\to1$; and
\item $\displaystyle\limsup_{n\to\infty}
 \norm{\Delta_n^*a\Delta_n}\leq |\psi(a)|$ for every $a\in A$.
\end{enumerate}
Clou\^atre asks in Question 1 of the section ``Characteristic sequences'':

\begin{quote}
If $\psi$ is a pure state on $A$ with the unique extension property with
respect to $S$, must there always exist a characteristic sequence in $S$ for
$\psi$?
\end{quote}

The original question and its surrounding discussion appear on PDF page 18
of the source and are reproduced in Figure~\ref{fig:source}.

\begin{figure}[ht]
\centering
\includegraphics[width=\textwidth]{figures/open_problem_crop.png}
\caption{Question 1 in arXiv:1709.01649v2, PDF page 18.}
\label{fig:source}
\end{figure}

Here the unique extension property (UEP) for a state means that it is the only
state on $A$ agreeing with its restriction to $S$.  Every pure state of
$\Md$ has the form
\[
 \psi(a)=\langle a\xi,\xi\rangle=\Tr(pa),
 \qquad p=\xi\xi^*,\quad \norm{\xi}=1,
\]
where $p$ is its rank-one support projection.

\section{Proof intuition}

In finite dimension, the asymptotic object in the question cannot hide at
infinity.  A characteristic sequence has a norm-convergent subsequence.  Its
limit $\Delta$ must be killed on the left by $1-p$, because condition (c)
applied to $a=1-p$ has right-hand side zero.  Hence $\Delta$ has range in the
one-dimensional range of $p$.  The norm and state normalization then force
$\Delta=p$.  Consequently, a characteristic sequence inside $S$ exists if
and only if $p\in S$.

The problem is therefore to make $p$ uniquely determined among density
matrices by its values on $S$, without putting $p$ itself in $S$.  In $M_3$
this is possible with a trace-orthogonal hyperplane.  Let
\[
 H=\operatorname{diag}(1,1,-2),\qquad
 S=\{a\in M_3:\Tr(Ha)=0\},\qquad p=E_{11}.
\]
Every density matrix agreeing with $p$ on $S$ lies on the affine line
$p+tH$.  Positivity of its second and third diagonal entries requires both
$t\geq0$ and $t\leq0$, so $t=0$.  Thus UEP holds, but $p\notin S$.

In $M_2$ there is no room for this two-sided positivity trap.  If $p\notin S$,
a trace-orthogonal self-adjoint direction through $p$ always points into the
state space for one sufficiently small choice of sign.  It produces a second
state with the same restriction, contradicting UEP.  This supplies the sharp
dimension threshold.

\section{The support-projection criterion}

\begin{lemma}\label{lem:criterion}
Let $S\subset\Md$ be an operator system, let
$\psi(a)=\langle a\xi,\xi\rangle$ be a pure state, and let
$p=\xi\xi^*$.  Then $\psi$ admits a characteristic sequence contained in
$S$ if and only if $p\in S$.
\end{lemma}

\begin{proof}
If $p\in S$, take $\Delta_n=p$ for all $n$.  Then $\norm p=1$,
$\psi(p)=1$, and
\[
 p a p=\psi(a)p,
 \qquad \norm{pap}=|\psi(a)|,
\]
so all three defining conditions hold.

Conversely, suppose $(\Delta_n)\subset S$ is characteristic.  The unit ball
of the finite-dimensional space $S$ is compact, so some subsequence converges
in norm to an element $\Delta\in S$.  Conditions (a) and (b) give
$\norm\Delta=1$ and $\psi(\Delta)=1$.  Put $q=1-p$.  Since $\psi(q)=0$,
condition (c), followed along the convergent subsequence, yields
\[
 \norm{\Delta^*q\Delta}=0.
\]
Thus $q\Delta=0$, or equivalently $p\Delta=\Delta$.  Moreover
$p\Delta p=\psi(\Delta)p=p$, and hence $\Delta p=p$.  Put
$x=\Delta q$.  Then $\Delta=p+x$, the initial spaces of $p$ and $x$ are
orthogonal, and the range of $x$ lies in the one-dimensional range of $p$.
Consequently
\[
 \Delta\Delta^*=p+xx^*=(1+\norm{x}^2)p.
\]
Since $\norm\Delta=1$, this forces $x=0$.  Hence $\Delta=p\in S$.
\end{proof}

\section{The counterexample and the sharp threshold}

\begin{theorem}\label{thm:main}
The source question has a negative answer already for $A=M_3(\mathbb C)$.
More precisely:
\begin{enumerate}[label=(\roman*)]
\item for every $d\geq3$, there is a generating operator system
$S\subset\Md$ and a pure state $\psi$ having the UEP with respect to $S$ but
no characteristic sequence in $S$;
\item for $d=1,2$, every pure state on $\Md$ having the UEP with respect to
an operator system $S$ does admit a characteristic sequence in $S$.
\end{enumerate}
Thus $d=3$ is the least full-matrix dimension in which the implication fails.
\end{theorem}

\begin{proof}
Fix $d\geq3$ and set
\[
 p=E_{11},\qquad H=\operatorname{diag}(1,1,-2,0,\ldots,0),
 \qquad S=\{a\in\Md:\Tr(Ha)=0\}.
\]
Because $H=H^*$ and $\Tr H=0$, the complex hyperplane $S$ is self-adjoint
and contains the identity; hence it is an operator system.  It contains every
off-diagonal matrix unit $E_{ij}$.  Products $E_{ij}E_{ji}=E_{ii}$ then show
that $C^*(S)=\Md$.  Also
\[
 \Tr(Hp)=1,
\]
so $p\notin S$.

Let $\psi(a)=\Tr(pa)=a_{11}$.  This is a pure state.  Suppose a state
$\varphi(a)=\Tr(\rho a)$ agrees with $\psi$ on $S$, where
$\rho\geq0$ and $\Tr\rho=1$.  The trace pairing on $\Md$ is nondegenerate,
and $S$ is the kernel of the nonzero functional $a\mapsto\Tr(Ha)$.
Therefore its annihilator is $\mathbb C H$, and agreement on $S$ gives
\[
 \rho-p=tH
\]
for some scalar $t$.  Self-adjointness makes $t$ real.  The second and third
diagonal entries of $\rho=p+tH$ are respectively $t$ and $-2t$.  Positivity
forces both to be nonnegative, hence $t=0$.  Thus $\rho=p$, proving that
$\psi$ has the UEP.  Since $p\notin S$, Lemma~\ref{lem:criterion} proves that
$\psi$ has no characteristic sequence in $S$.  This proves (i), including the
explicit $M_3$ counterexample.

For (ii), the $d=1$ case is immediate.  Let $d=2$, let $p$ be the rank-one
support of a pure state $\psi$, and suppose that $p\notin S$.  Regard
$M_2(\mathbb C)_{\rm sa}$ as a real inner-product space under
$(x,y)\mapsto\Tr(xy)$.  Since $p\notin S_{\rm sa}$, finite-dimensional
separation gives a self-adjoint $H$ such that
\[
 \Tr(Hs)=0\quad(s\in S),\qquad a:=\Tr(Hp)\neq0.
\]
The first equality holds for all complex $s\in S$ by decomposing $s$ into
self-adjoint real and imaginary parts.  Since $1\in S$, we also have
$\Tr H=0$.  In a basis with $p=E_{11}$, write
\[
 H=\begin{pmatrix}a&z\\ \overline z&-a\end{pmatrix}.
\]
Choose a nonzero real $t$ of sign opposite to $a$ and with magnitude so small
that
\[
 |t|<\min\left\{\frac1{|a|},
 \frac{|a|}{a^2+|z|^2}\right\}.
\]
Then $\rho=p+tH$ has trace one, both diagonal entries are positive, and
\[
 \det\rho=-ta-t^2(a^2+|z|^2)>0.
\]
Hence $\rho$ is a positive definite density matrix distinct from $p$.
Moreover $\Tr((\rho-p)s)=t\Tr(Hs)=0$ for all $s\in S$, so $\rho$ defines a
second state agreeing with $\psi$ on $S$.  Thus $p\notin S$ contradicts the
UEP.  We conclude that UEP forces $p\in S$ in $M_2$, and the constant sequence
$\Delta_n=p$ is characteristic by Lemma~\ref{lem:criterion}.
\end{proof}

\section{Verification, scope, and novelty}

The proof reduces every asymptotic statement to finite-dimensional norm
compactness.  The $M_d$ UEP calculation uses all states, not merely pure
states: every state is represented by a density matrix, and the annihilator
of the displayed hyperplane is exactly one-dimensional.  The example's
operator system generates the ambient matrix algebra, so the failure is not
caused by placing $S$ inside a proper $C^*$-subalgebra.

The theorem is a full counterexample to the unrestricted yes/no question and
a complete classification by dimension for full matrix algebras.  It does not
classify arbitrary finite-dimensional direct sums or infinite-dimensional
$C^*$-algebras, which is unnecessary for the negative answer.

A bounded novelty sweep searched the run's lightweight indexes and primary
arXiv sources for the exact question, characteristic sequences, UEP, and
pure-state restrictions.  It found the source paper and Clou\^atre's later
paper on restrictions of pure states and excision~\cite{ClouatreRestrictions},
but no explicit answer to Question 1 or the matrix construction above.  A
different existing run packet associated with arXiv:2306.02365 concerns
failure of a distinct pure-extension property and is not the present result.
Novelty confidence is therefore moderate rather than an assertion of
exhaustive priority.  No external theorem is used in the proof.

\paragraph{Human review checklist.}
Check the limiting argument in Lemma~\ref{lem:criterion}; verify the
one-dimensional trace annihilator and the two diagonal positivity constraints
in the $d\geq3$ example; and verify the determinant sign and small-$t$ bounds
in the $M_2$ converse.  These are the only substantive proof points.

\begin{thebibliography}{9}
\bibitem{ClouatreSource}
R. Clou\^atre,
\emph{Non-commutative peaking phenomena and a local version of the
hyperrigidity conjecture}, arXiv:1709.01649v2 (2018), especially Question 1,
PDF p.~18.

\bibitem{ClouatreRestrictions}
R. Clou\^atre,
\emph{Restrictions of pure states to subspaces of $C^*$-algebras},
arXiv:2306.02365 (2023; later revisions).
\end{thebibliography}

\end{document}
